\section*{Exercises}
\addcontentsline{toc}{section}{Exercises}

\subsection*{A. Theory}

\begin{exercise}
\exerciseid{11.A1}
Prove that if $Av=\lambda v$ and $v\neq 0$, then $A^kv=\lambda^kv$ for every
integer $k\geq 0$.
\end{exercise}

\begin{exercise}
\exerciseid{11.A2}
Prove that eigenvectors belonging to two distinct eigenvalues are linearly
independent.
\end{exercise}

\subsection*{B. By Hand}

\begin{exercise}
\exerciseid{11.B1}
Find the eigenvalues and eigenvectors of
\[
  \begin{bmatrix}2&0\\0&3\end{bmatrix}.
\]
\end{exercise}

\begin{exercise}
\exerciseid{11.B2}
For
\[
  A=\begin{bmatrix}1&2\\0&3\end{bmatrix},
\]
verify by hand that $(1,0)^T$ and $(1,1)^T$ are eigenvectors. Build $P$ from
these eigenvectors and compute $P^{-1}AP$.
\end{exercise}

\subsection*{C. Python / NumPy}

\begin{exercise}
\exerciseid{11.C1}
Use \texttt{np.linalg.eig} on a diagonal matrix and on a non-diagonal matrix
with known eigenvectors. Verify $Av=\lambda v$ for each eigenpair.
\end{exercise}

\begin{exercise}
\exerciseid{11.C2}
Choose a diagonalizable $2\times 2$ matrix $A$ and an initial vector $x_0$.
Compute $x_{t+1}=Ax_t$ for $t=0,\ldots,10$ and describe which eigenvector
direction dominates.
\end{exercise}

\subsection*{D. Challenge}

\begin{exercise}
\exerciseid{11.D1}
Give a $2\times 2$ matrix with only one eigenvector direction. Explain why it
cannot be diagonalized.
\end{exercise}
